\documentclass[11pt,oneside]{article}
\usepackage[utf8]{inputenc}
\usepackage[english,russian]{babel}
\usepackage{xcolor}
\usepackage{hyperref}

\usepackage{graphicx}
\usepackage{wrapfig}
\graphicspath{ {imgs/} }

\newcommand{\hhref}[2]{\href{#1}{\color{blue}#2}}
 
\begin{document}


\begin{center}
	{\huge Polyarniy Nikolay}
\end{center}


\vspace{-9pt}
\section*{\textbf{Work Experience}}
\vspace{-9pt}

\begin{description}
  \item[ - Agisoft] \hfill \\
    \textbf{Since April 2016}
  
    \textbf{Mathematician-Programmer}
    
    Metashape developer (former \hhref{https://en.wikipedia.org/wiki/PhotoScan}{PhotoScan}, see \hhref{http://www.agisoft.com/community/showcase/}{showcase} and \hhref{http://www.agisoft.com/community/blog/}{blog}).
    Mentor of students internship.
    Currently focused on detailed, fast and scalable polygonal model reconstruction.

    Computer vision, OpenCL and CUDA, computational geometry.
  \item[ - Transas] \hfill \\
    \textbf{October 2014 - March 2016}

    \textbf{Mathematician-Programmer}

    Developed server producing 3D landscape reconstruction and true orthophoto stitching from photos taken by UAV (\hhref{http://polarnick239.github.io/old/cv/Monoceros1.pdf}{presentation}, \hhref{http://polarnick239.github.io/old/cv/Monoceros2.pdf}{second presentation}).

    OpenCV, OpenCL, Python, Cython, Ceres-solver.
  \item[ - Yandex.Money] \hfill \\
    \textbf{February 2014 – October 2014:} Software Developer (Java backend)
  \item[ - DevExperts] \hfill \\
    \textbf{April 2013 – September 2013:} Software Developer (Java backend)

\end{description}


\vspace{-9pt}
\section*{\textbf{Skills}}
\vspace{-9pt}

\begin{itemize}
    \item{\textbf{Computer Vision}}: structure from motion, multiple view geometry, magic (like  \hhref{http://www.agisoft.com/index.php?id=49}{masks transfering}).
    Better than state of the art results in 3D polygonal surface reconstruction.
    Fast and occlusion-aware depth maps reconstruction.
    Dense point cloud classification: ground/road/building/vegetation/car.

    \item{\textbf{Computational geometry, CGAL}}: computations with absolute accuracy, algorithms and structures like Delaunay triangulation

    \item{\textbf{OpenCL, CUDA, OpenGL, WebGL}}: GPGPU computations, modern shaders, ray tracing. Can configure: GPU-passthrough in KVM, OpenGL in headless servers like EC2 g2/g3/p2/p3, simultaneous usage of NVIDIA+AMD.

    \item{\textbf{C++, Python, Java}}
\end{itemize}


\vspace{-9pt}
\section*{\textbf{Activities}}
\vspace{-9pt}

\begin{itemize}
    \item{\textbf{GPGPU course}}: developed GPGPU OpenCL \hhref{https://compscicenter.ru/courses/video_cards_computation/2018-autumn/classes/}{course} in CSC. Video recordings available on \hhref{https://www.youtube.com/watch?v=L79PgDOcVfw&list=PLlb7e2G7aSpTgwAm0GBkvn5XA0NokovJJ}{youtube}, tasks available on \hhref{https://github.com/GPGPUCourse2018}{github}.

    \item{\textbf{Open-source}: OpenCL \hhref{https://github.com/PolarNick239/OpenMeanShift}{modification} of EDISON mean shift segmentation. \hhref{https://github.com/opencv/opencv/pull/6078}{Implemented} Python bindings for OpenCL algorithms in OpenCV. Contributions to OpenCV, PyOpenCL, jupyter qtconsole and others. GPU monitoring in i3pystatus. Was a developer of \hhref{http://jtalks.org/}{jtalks.org} project (administration control component)}.

    \item{\textbf{Hackathons}}: four awards on hackatons. Third place on \hhref{https://career.luxoft.com/lp/hack-cv/}{HackCV} (traffic signs recognition), \hhref{http://hackday.ru/sciencehackday-2/projects\#project-1400}{Science Hackday \#2} (Startup nomination), \hhref{http://hackday.ru/hackday-36/projects\#project-1121}{Hackday\#36} (Autodesk 3D-web nomination), \hhref{https://www.hackerleague.org/hackathons/jetbrains-edtech-hackathon/blogposts/53655896e24d32cfbd000006}{HackEdu} by JetBrains (third place). Participation in \hhref{http://www.hackjunction.com/}{Junction 2016, 2017}.

    \item{\textbf{Conferences}}: participated in \hhref{http://3dv18.uniud.it/}{3DV 2018}, will participate in \hhref{http://www.3d-arch.org/}{3D-ARCH 2019}.

    \item{\textbf{Magister Ludi}}: \hhref{http://239.ru}{PML №239} programming teacher. \hhref{http://polarnick.com/blogs/239/2018/index_ru.html}{10-5} since september 2018. \hhref{http://polarnick239.github.io/239/201612}{9-1, 10-1, 11-1} in winter 2016/2017.

    \item{\textbf{Mathematical circle}}: PML №239 mathematical circle teacher, SPb school team mentor, mathematical contests volunteer.

    \item{\textbf{Programming contests}}: contests organization and participation.
\end{itemize}

\vspace{-9pt}
\section*{\textbf{Education}}
\vspace{-9pt}

\begin{itemize}
    \item{Computer Science Center}
    \item{ITMO University, Computer Technologies}
    \item{PML №239, mathematical circle, programming contests}
\end{itemize}

\begin{wrapfigure}{r}{0.22\textwidth}
    \centering
    \includegraphics[width=0.22\textwidth]{unicorn.png}
\end{wrapfigure}

\vspace{-9pt}
\section*{\textbf{Contacts}}
\vspace{-9pt}

\begin{itemize}

    \item{\textbf{\hhref{mailto:PolarNick239@gmail.com}{PolarNick239@gmail.com}}}

    \item{\textbf{\hhref{http://polarnick239.github.io/index_ru.html}{PolarNick.ru}}}

    \item{\textbf{\hhref{https://github.com/PolarNick239}{GitHub}}}
    
    \item{\textbf{\hhref{https://www.linkedin.com/in/nickolay-polyarniy-61393b7b}{LinkedIn}}}
    
    
    
\end{itemize}

Last updated: 20.01.2019

\end{document}
